%%%%%%%%%%%%%%%%%%%%%%%%%%%%%%%%%%%%%%%%%%%%%%%%%%%%%%%%%%%%%%%%%%%
% Header: General study / work / project template
%%%%%%%%%%%%%%%%%%%%%%%%%%%%%%%%%%%%%%%%%%%%%%%%%%%%%%%%%%%%%%%%%%%

\documentclass[11pt,a4paper]{scrartcl}

% Layout
\usepackage[margin=2.5cm]{geometry}
\usepackage{setspace}
\setstretch{1.15}

% General packages
\usepackage{graphicx}
\usepackage{enumitem}
\usepackage[most]{tcolorbox}
\tcbuselibrary{breakable}
\usepackage[breaklinks=true,colorlinks=true,linkcolor=blue,urlcolor=blue,citecolor=blue]{hyperref}
\usepackage{amsmath,amsthm,amssymb,mathtools}
\usepackage{braket}
\usepackage{icomma}
\usepackage[T1]{fontenc}
\usepackage[utf8]{inputenc}
\usepackage[ngerman]{babel}

% Units / tables / floats / references
\usepackage[locale = DE,space-before-unit=true,per-mode = symbol]{siunitx}
\usepackage{booktabs,multirow}
\usepackage{placeins}
\usepackage[natbib,abbreviate=true,doi=false,style=numeric-comp,giveninits=true,sorting=none]{biblatex}
\usepackage{csquotes}

% Optional math shortcuts
\newcommand{\R}{\mathbb{R}}
\newcommand{\N}{\mathbb{N}}
\newcommand{\Q}{\mathbb{Q}}
\newcommand{\abs}[1]{\left| #1 \right|}
\newcommand{\norm}[1]{\left\lVert #1 \right\rVert}
\newcommand{\ol}[1]{\overline{#1}}
\newcommand{\op}[1]{\hat{#1}}
\newcommand{\e}{\mathrm e}
\newcommand{\ii}{\mathrm i}
\newcommand{\eps}{\varepsilon}
\newcommand{\Ylm}{Y_l^m}
\newcommand{\Pex}{\pi_{12}}

% Box styles
\newtcolorbox{calculationbox}[1][]{
    title=Calculation,
    breakable,
    enhanced,
    colback=white,
    colframe=black!65,
    #1
}

\newtcolorbox{solutionbox}[1][]{
    title=Solution,
    breakable,
    enhanced,
    colback=white,
    colframe=black!65,
    #1
}

\newtcolorbox{answerbox}[1][]{
    title=Answer,
    breakable,
    enhanced,
    colback=white,
    colframe=black!65,
    #1
}

\newtcolorbox{notebox}[1][]{
    title=Notes,
    breakable,
    enhanced,
    colback=white,
    colframe=black!65,
    #1
}

\newtcolorbox{sidecalcbox}[1][]{
    title=Side Calculation,
    breakable,
    enhanced,
    colback=white,
    colframe=black!65,
    #1
}

%%%%%%%%%%%%%%%%%%%%%%%%%%%%%%%%%%%%%%%%%%%%%%%%%%%%%%%%%%%%%%%%%%%
\begin{document}

    \title{Theoretische Physik 2 -- Blatt 10}
    \author{Tobias Laser}
    \date{\today}
    \maketitle
    \vfill
    \thispagestyle{empty}
    \newpage

    \pagenumbering{arabic}

    % ================================================================
    \paragraph{23. Identische Teilchen im Kasten:}

    Zwei identische Spin-$\frac12$-Teilchen bewegen sich in einem eindimensionalen Kastenpotential
    \begin{align*}
        V(x)=
        \begin{cases}
            \infty, & x<0 \text{ oder } x>L,\\
            0, & 0\le x\le L.
        \end{cases}
    \end{align*}
    Die Einteilchen-Eigenfunktionen und Einteilchen-Energien im Kasten sind
    \begin{align*}
        \varphi_n(x)=\sqrt{\frac{2}{L}}\sin\left(\frac{n\pi x}{L}\right),
        \qquad
        \eps_n=\frac{\hbar^2\pi^2 n^2}{2mL^2},
        \qquad n=1,2,3,\dots
    \end{align*}
    Da die Teilchen Spin-$\frac12$-Fermionen sind, muss die volle Wellenfunktion unter Teilchenaustausch antisymmetrisch sein.

    \begin{enumerate}[label=(\alph*)]
        \item Finden Sie den Grundzustand und die Grundzustandsenergie für den Fall, in dem sich die Teilchen in einem Triplettzustand befinden.

        \begin{calculationbox}
            Ein Triplettzustand ist im Spinanteil symmetrisch. Damit die gesamte Wellenfunktion für zwei identische Fermionen antisymmetrisch ist, muss der Ortsanteil antisymmetrisch sein.

            \begin{calculationbox}[title=Symmetriebedingung]
                Für Fermionen gilt
                \begin{align*}
                    \Pex\Psi(x_1,x_2;\sigma_1,\sigma_2)=-\Psi(x_1,x_2;\sigma_1,\sigma_2).
                \end{align*}
                Im Triplett ist der Spinanteil symmetrisch:
                \begin{align*}
                    \Pex\chi_{S=1,M}=+\chi_{S=1,M}.
                \end{align*}
                Also muss für den Ortsanteil gelten
                \begin{align*}
                    \Pex\psi_{\mathrm{ort}}(x_1,x_2)=-\psi_{\mathrm{ort}}(x_1,x_2).
                \end{align*}
            \end{calculationbox}

            \begin{calculationbox}[title=Kleinste antisymmetrische Ortswellenfunktion]
                Eine antisymmetrische Ortswellenfunktion kann nicht beide Teilchen im selben Einteilchenzustand haben, denn
                \begin{align*}
                    \frac{1}{\sqrt2}\left[\varphi_1(x_1)\varphi_1(x_2)-\varphi_1(x_2)\varphi_1(x_1)\right]=0.
                \end{align*}
                Deshalb müssen die beiden niedrigsten verschiedenen Einteilchenzustände besetzt werden, also $n=1$ und $n=2$:
                \begin{align*}
                    \psi_A(x_1,x_2)
                    =
                    \frac{1}{\sqrt2}
                    \left[
                        \varphi_1(x_1)\varphi_2(x_2)
                        -
                        \varphi_2(x_1)\varphi_1(x_2)
                    \right].
                \end{align*}
            \end{calculationbox}

            \begin{calculationbox}[title=Grundzustandsenergie]
                Die Energie ist die Summe der besetzten Einteilchenenergien:
                \begin{align*}
                    E_T^{(0)}
                    &=\eps_1+\eps_2\\
                    &=\frac{\hbar^2\pi^2}{2mL^2}
                    +
                    \frac{4\hbar^2\pi^2}{2mL^2}\\
                    &=\frac{5\hbar^2\pi^2}{2mL^2}.
                \end{align*}

                \begin{solutionbox}[title=Triplett-Grundzustand]
                    \begin{align*}
                        \Psi_{T,M}(x_1,x_2)
                        =
                        \psi_A(x_1,x_2)\chi_{1,M},
                        \qquad
                        E_T^{(0)}=\frac{5\hbar^2\pi^2}{2mL^2}.
                    \end{align*}
                \end{solutionbox}
            \end{calculationbox}
        \end{calculationbox}

        \item Finden Sie den Grundzustand und die Grundzustandsenergie für den Fall, in dem sich die Teilchen in einem Singulettzustand befinden.

        \begin{calculationbox}
            Ein Singulettzustand ist im Spinanteil antisymmetrisch. Deshalb muss der Ortsanteil symmetrisch sein, damit die gesamte Wellenfunktion wieder antisymmetrisch ist.

            \begin{calculationbox}[title=Symmetriebedingung]
                Im Singulett gilt
                \begin{align*}
                    \Pex\chi_{0,0}=-\chi_{0,0}.
                \end{align*}
                Da die Gesamtwellenfunktion antisymmetrisch sein muss, folgt
                \begin{align*}
                    \Pex\psi_{\mathrm{ort}}(x_1,x_2)=+\psi_{\mathrm{ort}}(x_1,x_2).
                \end{align*}
            \end{calculationbox}

            \begin{calculationbox}[title=Kleinste symmetrische Ortswellenfunktion]
                Für den symmetrischen Ortsanteil dürfen beide Fermionen denselben Einteilchen-Ortszustand besetzen, weil der Spinanteil schon antisymmetrisch ist. Der niedrigste Zustand ist daher
                \begin{align*}
                    \psi_S(x_1,x_2)
                    =
                    \varphi_1(x_1)\varphi_1(x_2).
                \end{align*}
            \end{calculationbox}

            \begin{calculationbox}[title=Grundzustandsenergie]
                Die Energie ist
                \begin{align*}
                    E_S^{(0)}
                    &=\eps_1+\eps_1\\
                    &=2\frac{\hbar^2\pi^2}{2mL^2}\\
                    &=\frac{\hbar^2\pi^2}{mL^2}.
                \end{align*}

                \begin{solutionbox}[title=Singulett-Grundzustand]
                    \begin{align*}
                        \Psi_S(x_1,x_2)
                        =
                        \varphi_1(x_1)\varphi_1(x_2)\chi_{0,0},
                        \qquad
                        E_S^{(0)}=\frac{\hbar^2\pi^2}{mL^2}.
                    \end{align*}
                \end{solutionbox}
            \end{calculationbox}
        \end{calculationbox}

        \item Wir nehmen an, die Teilchen unterliegen einer sehr kurzreichweitigen attraktiven Wechselwirkung
        \begin{align*}
            \op V=-\lambda\delta(x_1-x_2),
            \qquad \lambda>0.
        \end{align*}
        Diskutieren Sie mittels Störungstheorie den Effekt dieses Potentials auf die Grundzustände.

        \begin{calculationbox}
            In erster Ordnung Störungstheorie ist die Energieverschiebung
            \begin{align*}
                \Delta E^{(1)}
                =
                \bra{\psi}\op V\ket{\psi}
                =
                -\lambda
                \int_0^L\int_0^L
                \abs{\psi(x_1,x_2)}^2\delta(x_1-x_2)\,dx_1\,dx_2.
            \end{align*}
            Mit der Deltafunktion wird daraus
            \begin{align*}
                \Delta E^{(1)}
                =
                -\lambda\int_0^L\abs{\psi(x,x)}^2\,dx.
            \end{align*}

            \begin{calculationbox}[title=Triplett-Ortszustand]
                Für den antisymmetrischen Ortszustand gilt am Ort $x_1=x_2=x$
                \begin{align*}
                    \psi_A(x,x)
                    &=
                    \frac{1}{\sqrt2}
                    \left[
                        \varphi_1(x)\varphi_2(x)
                        -
                        \varphi_2(x)\varphi_1(x)
                    \right]\\
                    &=0.
                \end{align*}
                Daher ist die erste Ordnungsverschiebung
                \begin{align*}
                    \Delta E_T^{(1)}=0.
                \end{align*}

                \begin{solutionbox}[title=Effekt im Triplett]
                    \begin{align*}
                        \Delta E_T^{(1)}=0.
                    \end{align*}
                    Die kurzreichweitige Wechselwirkung wirkt nicht in erster Ordnung, weil die antisymmetrische Ortswellenfunktion bei Teilchenkoinzidenz verschwindet.
                \end{solutionbox}
            \end{calculationbox}

            \begin{calculationbox}[title=Singulett-Ortszustand]
                Im Singulett ist
                \begin{align*}
                    \psi_S(x_1,x_2)=\varphi_1(x_1)\varphi_1(x_2),
                    \qquad
                    \psi_S(x,x)=\varphi_1(x)^2.
                \end{align*}
                Damit
                \begin{align*}
                    \Delta E_S^{(1)}
                    &=-\lambda\int_0^L\abs{\varphi_1(x)}^4\,dx\\
                    &=-\lambda\int_0^L
                    \left(\frac{2}{L}\sin^2\left(\frac{\pi x}{L}\right)\right)^2dx\\
                    &=-\lambda\frac{4}{L^2}\int_0^L\sin^4\left(\frac{\pi x}{L}\right)dx.
                \end{align*}
                Mit
                \begin{align*}
                    \int_0^L\sin^4\left(\frac{\pi x}{L}\right)dx=\frac{3L}{8}
                \end{align*}
                folgt
                \begin{align*}
                    \Delta E_S^{(1)}
                    =
                    -\lambda\frac{4}{L^2}\frac{3L}{8}
                    =
                    -\frac{3\lambda}{2L}.
                \end{align*}

                \begin{solutionbox}[title=Effekt im Singulett]
                    \begin{align*}
                        \Delta E_S^{(1)}=-\frac{3\lambda}{2L},
                        \qquad
                        E_S\approx \frac{\hbar^2\pi^2}{mL^2}-\frac{3\lambda}{2L}.
                    \end{align*}
                    Die attraktive Delta-Wechselwirkung senkt den Singulett-Grundzustand in erster Ordnung ab.
                \end{solutionbox}
            \end{calculationbox}
        \end{calculationbox}
    \end{enumerate}

    % ================================================================
    \paragraph{24. Relativer Drehimpuls identischer Teilchen:}

    Wir betrachten zwei identische Teilchen mit Wechselwirkungspotential, das nur vom Betrag des relativen Abstands abhängt:
    \begin{align*}
        \op H
        =
        \frac{\vec p_1^{\,2}}{2m}
        +
        \frac{\vec p_2^{\,2}}{2m}
        +
        V(\abs{\vec r_1-\vec r_2}).
    \end{align*}

    \begin{enumerate}[label=(\alph*)]
        \item Geben Sie $\op H$ in den Schwerpunkts- und Relativkoordinaten
        \begin{align*}
            \vec R=\frac12(\vec r_1+\vec r_2),
            \qquad
            \vec P=\vec p_1+\vec p_2,
            \qquad
            \vec r=\vec r_1-\vec r_2,
            \qquad
            \vec p=\frac12(\vec p_1-\vec p_2)
        \end{align*}
        an und zeigen Sie, dass $[H,\vec P]=[H,\vec L^2]=[H,L_z]=0$ mit $\vec L=\vec r\times\vec p$. Zeigen Sie außerdem die Form der Eigenfunktionen und geben Sie die Radialgleichung für $R_{n,l}(r)$ an.

        \begin{calculationbox}
            Wir lösen zuerst die Impulse nach $\vec p_1$ und $\vec p_2$ auf.

            \begin{calculationbox}[title=Transformation der kinetischen Energie]
                Aus
                \begin{align*}
                    \vec P=\vec p_1+\vec p_2,
                    \qquad
                    \vec p=\frac12(\vec p_1-\vec p_2)
                \end{align*}
                folgt
                \begin{align*}
                    \vec p_1=\frac{\vec P}{2}+\vec p,
                    \qquad
                    \vec p_2=\frac{\vec P}{2}-\vec p.
                \end{align*}
                Daher ist
                \begin{align*}
                    \frac{\vec p_1^{\,2}}{2m}+\frac{\vec p_2^{\,2}}{2m}
                    &=
                    \frac{1}{2m}
                    \left[
                        \left(\frac{\vec P}{2}+\vec p\right)^2
                        +
                        \left(\frac{\vec P}{2}-\vec p\right)^2
                    \right]\\
                    &=
                    \frac{1}{2m}
                    \left[
                        \frac{\vec P^{\,2}}{2}+2\vec p^{\,2}
                    \right]\\
                    &=
                    \frac{\vec P^{\,2}}{4m}+\frac{\vec p^{\,2}}{m}.
                \end{align*}

                \begin{solutionbox}[title=Hamiltonoperator in Schwerpunkts- und Relativkoordinaten]
                    \begin{align*}
                        \op H
                        =
                        \frac{\vec P^{\,2}}{4m}
                        +
                        \frac{\vec p^{\,2}}{m}
                        +
                        V(r),
                        \qquad r=\abs{\vec r}.
                    \end{align*}
                \end{solutionbox}
            \end{calculationbox}

            \begin{calculationbox}[title=Kommutatoren]
                Der Hamiltonoperator hängt von $\vec R$ nicht ab. Deshalb ist der Schwerpunktsimpuls erhalten:
                \begin{align*}
                    [\op H,\vec P]=0.
                \end{align*}
                Der relative Anteil
                \begin{align*}
                    \op H_{\mathrm{rel}}=\frac{\vec p^{\,2}}{m}+V(r)
                \end{align*}
                ist rotationsinvariant, weil $\vec p^{\,2}$ ein Skalar ist und $V$ nur von $r$ abhängt. Deshalb kommutiert $\op H_{\mathrm{rel}}$ mit dem relativen Drehimpuls:
                \begin{align*}
                    [\op H,\vec L^{\,2}]=0,
                    \qquad
                    [\op H,L_z]=0.
                \end{align*}
                Man kann also gemeinsame Eigenzustände von $\op H$, $\vec P$, $\vec L^{\,2}$ und $L_z$ wählen.
            \end{calculationbox}

            \begin{calculationbox}[title=Separationsansatz]
                Da der Schwerpunkt frei ist, ist der Schwerpunktanteil eine ebene Welle. Da der relative Hamiltonoperator zentral ist, trennt man den Winkelanteil in Kugelflächenfunktionen:
                \begin{align*}
                    \psi_{\vec P,n,l,m}(\vec R,\vec r)
                    =
                    \frac{1}{(2\pi\hbar)^{3/2}}
                    \e^{\ii\vec P\cdot\vec R/\hbar}
                    R_{n,l}(r)Y_l^m(\theta,\varphi).
                \end{align*}
            \end{calculationbox}

            \begin{calculationbox}[title=Radialgleichung]
                In Ortsdarstellung gilt für den relativen Impuls
                \begin{align*}
                    \vec p^{\,2}=-\hbar^2\nabla_r^2.
                \end{align*}
                Für Funktionen der Form $R_{n,l}(r)Y_l^m(\theta,\varphi)$ ist
                \begin{align*}
                    \nabla_r^2
                    =
                    \frac{1}{r^2}\frac{d}{dr}\left(r^2\frac{d}{dr}\right)
                    -
                    \frac{\vec L^{\,2}}{\hbar^2r^2},
                    \qquad
                    \vec L^{\,2}Y_l^m=\hbar^2l(l+1)Y_l^m.
                \end{align*}
                Wenn
                \begin{align*}
                    E=\frac{\vec P^{\,2}}{4m}+\eps_{n,l},
                \end{align*}
                dann erfüllt $R_{n,l}(r)$ die Gleichung
                \begin{align*}
                    \left[
                        -\frac{\hbar^2}{m}
                        \frac{1}{r^2}\frac{d}{dr}\left(r^2\frac{d}{dr}\right)
                        +
                        \frac{\hbar^2l(l+1)}{mr^2}
                        +
                        V(r)
                    \right]
                    R_{n,l}(r)
                    =
                    \eps_{n,l}R_{n,l}(r).
                \end{align*}

                \begin{solutionbox}[title=Eigenfunktionen und Radialgleichung]
                    \begin{align*}
                        \psi_{\vec P,n,l,m}(\vec R,\vec r)
                        =
                        \frac{1}{(2\pi\hbar)^{3/2}}
                        \e^{\ii\vec P\cdot\vec R/\hbar}
                        R_{n,l}(r)Y_l^m(\theta,\varphi),
                    \end{align*}
                    \begin{align*}
                        \left[
                            -\frac{\hbar^2}{m}
                            \frac{1}{r^2}\frac{d}{dr}\left(r^2\frac{d}{dr}\right)
                            +
                            \frac{\hbar^2l(l+1)}{mr^2}
                            +
                            V(r)
                        \right]
                        R_{n,l}(r)=\eps_{n,l}R_{n,l}(r).
                    \end{align*}
                \end{solutionbox}
            \end{calculationbox}
        \end{calculationbox}

        \item Die Teilchen seien Bosonen mit Spin $0$. Zeigen Sie, dass nur geradzahlige Werte für $l$ erlaubt sind.

        \begin{calculationbox}
            Für zwei identische Spin-$0$-Bosonen gibt es keinen nichttrivialen Spinanteil. Die Ortswellenfunktion selbst muss also symmetrisch unter Teilchenaustausch sein.

            \begin{calculationbox}[title=Wirkung des Teilchenaustauschs]
                Unter Austausch $1\leftrightarrow 2$ gilt
                \begin{align*}
                    \vec R\mapsto \vec R,
                    \qquad
                    \vec r=\vec r_1-\vec r_2\mapsto -\vec r.
                \end{align*}
                In Kugelkoordinaten bedeutet $\vec r\mapsto -\vec r$
                \begin{align*}
                    r\mapsto r,
                    \qquad
                    \theta\mapsto \pi-\theta,
                    \qquad
                    \varphi\mapsto \varphi+\pi.
                \end{align*}
                Der Schwerpunktfaktor und der Radialfaktor bleiben unverändert. Relevant ist also nur die Kugelflächenfunktion.
            \end{calculationbox}

            \begin{calculationbox}[title=Parität der Kugelflächenfunktionen]
                Mit der angegebenen Darstellung
                \begin{align*}
                    Y_l^m(\theta,\varphi)
                    =
                    N_{l,m}\e^{\ii m\varphi}(\sin\theta)^m
                    \frac{d^{l+m}}{d(\cos\theta)^{l+m}}(\sin\theta)^{2l}
                \end{align*}
                erhält man unter $\theta\mapsto \pi-\theta$ und $\varphi\mapsto \varphi+\pi$ den bekannten Paritätsfaktor
                \begin{align*}
                    Y_l^m(\pi-\theta,\varphi+\pi)=(-1)^lY_l^m(\theta,\varphi).
                \end{align*}
                Kurz gesagt: Kugelflächenfunktionen haben Parität $(-1)^l$.
            \end{calculationbox}

            \begin{calculationbox}[title=Bosonische Symmetriebedingung]
                Für Bosonen muss gelten
                \begin{align*}
                    \Pex\psi_{\vec P,n,l,m}(\vec R,\vec r)=+\psi_{\vec P,n,l,m}(\vec R,\vec r).
                \end{align*}
                Wegen $\vec r\mapsto -\vec r$ folgt
                \begin{align*}
                    (-1)^l=+1.
                \end{align*}

                \begin{solutionbox}[title=Erlaubte Werte für Spin-$0$-Bosonen]
                    \begin{align*}
                        l=0,2,4,\dots
                    \end{align*}
                    Für identische Spin-$0$-Bosonen sind nur gerade relative Drehimpulse erlaubt.
                \end{solutionbox}
            \end{calculationbox}
        \end{calculationbox}

        \item Die Teilchen seien Fermionen mit Spin $\frac12$. Welche Werte sind für $l$ erlaubt, wenn sich die Spins in einem Singulett- bzw. Triplettzustand befinden?

        \begin{calculationbox}
            Für zwei identische Fermionen muss die volle Wellenfunktion antisymmetrisch sein. Der Ortsanteil liefert unter Austausch den Faktor $(-1)^l$.

            \begin{calculationbox}[title=Spin-Symmetrie]
                Für zwei Spin-$\frac12$-Teilchen gilt:
                \begin{align*}
                    \text{Singulett }(S=0):&\quad \text{antisymmetrischer Spinanteil},\\
                    \text{Triplett }(S=1):&\quad \text{symmetrischer Spinanteil}.
                \end{align*}
                Also ist der Austauschfaktor des Spinanteils
                \begin{align*}
                    \eta_S=
                    \begin{cases}
                        -1, & S=0,\\
                        +1, & S=1.
                    \end{cases}
                \end{align*}
            \end{calculationbox}

            \begin{calculationbox}[title=Fermionische Antisymmetrie]
                Die Gesamtwellenfunktion transformiert mit
                \begin{align*}
                    \Pex\Psi=(-1)^l\eta_S\Psi.
                \end{align*}
                Für Fermionen muss $\Pex\Psi=-\Psi$ gelten. Daher
                \begin{align*}
                    (-1)^l\eta_S=-1.
                \end{align*}

                \begin{solutionbox}[title=Erlaubte relative Drehimpulse]
                    \begin{align*}
                        S=0\text{ Singulett:}\quad &(-1)^l(-1)=-1
                        \quad\Longrightarrow\quad l=0,2,4,\dots,\\
                        S=1\text{ Triplett:}\quad &(-1)^l(+1)=-1
                        \quad\Longrightarrow\quad l=1,3,5,\dots.
                    \end{align*}
                \end{solutionbox}
            \end{calculationbox}
        \end{calculationbox}

        \item Die Teilchen seien identische Spin-$1$-Bosonen. Welche zulässigen Werte für $l$ ergeben sich?

        \begin{calculationbox}
            Für zwei identische Spin-$1$-Bosonen muss die volle Wellenfunktion symmetrisch sein. Die möglichen Gesamtspins sind
            \begin{align*}
                S=0,1,2.
            \end{align*}
            Der Ortsanteil liefert wieder den Faktor $(-1)^l$.

            \begin{calculationbox}[title=Symmetrie des Spinanteils]
                Für zwei gleiche Spin-$1$-Teilchen gilt unter Teilchenaustausch
                \begin{align*}
                    \eta_S=(-1)^{2s-S}=(-1)^{2-S},
                    \qquad s=1.
                \end{align*}
                Damit
                \begin{align*}
                    S=0:&\quad \eta_S=+1,\\
                    S=1:&\quad \eta_S=-1,\\
                    S=2:&\quad \eta_S=+1.
                \end{align*}
                Also sind $S=0$ und $S=2$ symmetrische Spin-Zustände, während $S=1$ antisymmetrisch ist.
            \end{calculationbox}

            \begin{calculationbox}[title=Bosonische Symmetriebedingung]
                Für Bosonen gilt
                \begin{align*}
                    (-1)^l\eta_S=+1.
                \end{align*}

                \begin{solutionbox}[title=Erlaubte Werte für Spin-$1$-Bosonen]
                    \begin{align*}
                        S=0\text{ oder }S=2:
                        \quad l=0,2,4,\dots,\\
                        S=1:
                        \quad l=1,3,5,\dots.
                    \end{align*}
                \end{solutionbox}
            \end{calculationbox}
        \end{calculationbox}
    \end{enumerate}

    % ================================================================
    \paragraph{25. Bonusaufgabe: Spin-Bahn-Kopplung und Clebsch-Gordan-Koeffizienten:}

    Gesucht sind die Clebsch-Gordan-Koeffizienten für die Kopplung eines Bahndrehimpulses $l$ an den Elektronspin $s=\frac12$:
    \begin{align*}
        \ket{l,\tfrac12,J,M},
        \qquad
        J=l\pm\frac12.
    \end{align*}
    Für festes $M$ gibt es die beiden Produktbasiszustände
    \begin{align*}
        \ket a
        :=
        \ket{l,M+\tfrac12;\tfrac12,-\tfrac12},
        \qquad
        \ket b
        :=
        \ket{l,M-\tfrac12;\tfrac12,+\tfrac12}.
    \end{align*}

    \begin{calculationbox}
        Die beiden Zustände $\ket a$ und $\ket b$ haben beide Gesamtmagnetquantenzahl $M$. Wir diagonalisiseren $\vec J^{\,2}$ in diesem zweidimensionalen Raum.

        \begin{calculationbox}[title=Matrix von $\vec J^{\,2}$ in der Produktbasis]
            Es gilt
            \begin{align*}
                \vec J=\vec L+\vec S,
                \qquad
                \vec J^{\,2}=\vec L^{\,2}+\vec S^{\,2}+2L_zS_z+L_+S_-+L_-S_+.
            \end{align*}
            Auf der Basis $(\ket a,\ket b)$ erhält man
            \begin{align*}
                \frac{\vec J^{\,2}}{\hbar^2}
                =
                \begin{pmatrix}
                    l(l+1)+\frac14-M & C\\
                    C & l(l+1)+\frac14+M
                \end{pmatrix},
            \end{align*}
            wobei
            \begin{align*}
                C
                =
                \sqrt{l(l+1)-\left(M+\frac12\right)\left(M-\frac12\right)}
                =
                \sqrt{\left(l+\frac12\right)^2-M^2}.
            \end{align*}
        \end{calculationbox}

        \begin{calculationbox}[title=Eigenwerte]
            Die möglichen Gesamtbahndrehimpulse sind
            \begin{align*}
                J_+=l+\frac12,
                \qquad
                J_-=l-\frac12.
            \end{align*}
            Daher müssen die Eigenwerte von $\vec J^{\,2}$ gleich
            \begin{align*}
                \hbar^2J_+(J_++1)
                =
                \hbar^2\left(l+\frac12\right)\left(l+\frac32\right)
            \end{align*}
            beziehungsweise
            \begin{align*}
                \hbar^2J_-(J_-+1)
                =
                \hbar^2\left(l-\frac12\right)\left(l+\frac12\right)
            \end{align*}
            sein. Die zugehörigen normierten Eigenvektoren liefern die Clebsch-Gordan-Koeffizienten.
        \end{calculationbox}

        \begin{calculationbox}[title=Normierte Eigenzustände]
            Mit der Condon-Shortley-Phasenkonvention erhält man
            \begin{align*}
                \ket{l,\tfrac12,l+\tfrac12,M}
                &=
                \sqrt{\frac{l-M+\frac12}{2l+1}}
                \ket{l,M+\tfrac12;\tfrac12,-\tfrac12}
                +
                \sqrt{\frac{l+M+\frac12}{2l+1}}
                \ket{l,M-\tfrac12;\tfrac12,+\tfrac12},\\[0.7em]
                \ket{l,\tfrac12,l-\tfrac12,M}
                &=
                -\sqrt{\frac{l+M+\frac12}{2l+1}}
                \ket{l,M+\tfrac12;\tfrac12,-\tfrac12}
                +
                \sqrt{\frac{l-M+\frac12}{2l+1}}
                \ket{l,M-\tfrac12;\tfrac12,+\tfrac12}.
            \end{align*}
        \end{calculationbox}

        \begin{solutionbox}[title=Clebsch-Gordan-Koeffizienten]
            Für $J=l+\frac12$ gilt
            \begin{align*}
                \left\langle l,M+\frac12;\frac12,-\frac12\middle|l,\frac12,l+\frac12,M\right\rangle
                &=
                \sqrt{\frac{l-M+\frac12}{2l+1}},\\
                \left\langle l,M-\frac12;\frac12,+\frac12\middle|l,\frac12,l+\frac12,M\right\rangle
                &=
                \sqrt{\frac{l+M+\frac12}{2l+1}}.
            \end{align*}
            Für $J=l-\frac12$ gilt
            \begin{align*}
                \left\langle l,M+\frac12;\frac12,-\frac12\middle|l,\frac12,l-\frac12,M\right\rangle
                &=
                -\sqrt{\frac{l+M+\frac12}{2l+1}},\\
                \left\langle l,M-\frac12;\frac12,+\frac12\middle|l,\frac12,l-\frac12,M\right\rangle
                &=
                \sqrt{\frac{l-M+\frac12}{2l+1}}.
            \end{align*}
        \end{solutionbox}

        \begin{notebox}[title=Hinweis zur Phase]
            Das relative Minuszeichen im Zustand mit $J=l-\frac12$ ist eine Konvention. Physikalische Wahrscheinlichkeiten hängen von den Betragsquadraten der Clebsch-Gordan-Koeffizienten ab; die Standardtabellen verwenden typischerweise die oben gewählte Condon-Shortley-Konvention.
        \end{notebox}
    \end{calculationbox}

\end{document}
